\documentclass[../main.tex]{subfiles}

\begin{document}
\section{Random Facts}
\subsection{The Polynomial is irreducible}
\begin{lemma}
  $f_n$ is irreducible over $\Q$.
\end{lemma}
\begin{proof}
  [Sketch:]
  Assume $n > 2$ (otherwise it is trivial).
  The winding number of $g_n$ as $x$ traverses the circle of radius $1 + \epsilon$ is $n-1$ for sufficiently small $\epsilon > 0$.
  Indeed, for such $\epsilon$, 
  \[
    |2x^{n-1}| > |x^n| + 1 \geq |x^n + 1|
  \]
  since 
  \[
    \left( 1+\epsilon \right)^n - 2\left( 1+\epsilon \right)^{n-1} + 1 < 0
  \]
  since 
  \[
    g'\left( 1 \right) = n - 2\left( n - 1 \right) = 2 - n < 0
  \]
  Therefore, it has $n-1$ roots of absolute value at most $1$.
  One of these roots is $1$, and hence $f_n$ has $n-2$ roots of absolute value at most $1$.

  Moreover, the only root of absolute value exactly $1$ is $1$ itself.
  Indeed, it suffices to show that the only root of $\hat{g_n}\left( x \right):= g_n\left( 1/ x \right) x^n = x^n - 2x + 1$ of absolute value $1$ is $1$ itself (since the roots of $\hat{g_n}$ are the inverses of the roots of $g_n$).
  However, any algebraic number $\alpha$ of absolute value $1$ is a root of unity, and hence if $g_n\left( \alpha \right) = 0$ then $1 - 2\alpha$ is a root of unity as well, and in particular, by the inverse triangle inequality,
  \[
    1 = |2\alpha - 1| \geq 2 - 1 = 1 
  \]
  and hence $\alpha = 1$.
   
  Therefore, $f_n$ has $n-2$ roots of absolute value strictly less than $1$.
  In other words, only one of its roots has absolute value greater than or equal to $1$ (in fact greater).
  If we assume towards contradiction that $f_n$ is reducible over $\Q$, then since it is monic, it is reducible over $\Z$ as well, and hence one of its factors is a monic polynomial with integer coefficients whose roots are all of absolute value strictly less than $1$ which is a contradiction (because then for example the norm, which is the free coefficient up to a sign, is less than $1$ in absolute value).
\end{proof}


\subsection{The Action of the Inertia Group on the Roots of $f_n$}
\begin{lemma}
  Let $\P$ be a prime of $\O_{L_n}$ lying above $\left( 2 \right) \subseteq \Z$.
  Then the inertia group $I_{\P}$ acts transitively on each inertia block of $\P$.
\end{lemma}
\begin{proof}
  \editorial{Tweak the proof to apply to the tamely ramified case as well}

  Let $B$ be an inertia block of $\P$.
  It suffices to show that all elements of $B$ are conjugate over $\Q_2^\un$, the maximal unramified extension of $\Q_2$.
  To show that in turn, it suffices to show that the degree of one of these elements over $\Q_2^\un$ is at least the size of $B$.
  Indeed, assume that $\alpha \in B$ and assume that the degree of $\alpha$ over $\Q_2^\un$ is at least $e = \left| B \right|$.
  Assume that $\beta$ is a conjugate of $\alpha$ over $\Q_2^\un$.
  Then there is an automorphism $\sigma$ of $\Q_2^\alg$ over $\Q_2^\un$ such that $\sigma\left( \alpha \right) = \beta$.
  In particular, it reduces to a member of the inertia subgroup of $\Gal(L_n/\Q)$ at $\P$ that maps $\alpha$ to $\beta$ and hence $\alpha$ and $\beta$ are equivalent modulo $\P$ and hence $\beta \in B$.
  Therefore, for the number of conjugates of $\alpha$ over $\Q_2^\un$ to reach $e$, all elements of $B$ must be conjugate to $\alpha$ over $\Q_2^\un$.
  
  It remains to show that there is a root $\alpha \in B$ whose degree over $\Q_2^\un$ is at least $e$ or equivalently that there is $\beta \in \Q_2^\un\left( \alpha \right)$ whose $2$-adic valuation is in the interval $(0, 1/e]$.
  Indeed, let $k = n / 2^e$ and let $\beta$ be the unique primitive $k$-th root of unity whose image under $\pi_\P$ is the image of $B$ under $\pi_\P$.
  Then the free coefficient of $g_n\left( x + \beta \right)$ is 
  \[
    g_n\left( \beta \right) = 
    \beta^n - 2\beta^{n - 1} + 1 = 
    2\left( 1 - \beta^{n - 1} \right)
  \]
  whose $2$-adic valuation is exactly $1$.
  Therefore, since $g_n\left( x + \beta \right)$ has exactly $2^e$ roots whose $\P$-valuation is strictly positive, by Newton's polygon argument, there is a root $\alpha$ of $g_n\left( x + \beta \right)$ whose $\P$-valuation is not more than $1 / 2^e$.
\end{proof}



\subsection{The Discriminant of \( f_n \)}
\begin{lemma}
  \label{lem:2-adic-valuation-of-discriminant}
Assume that $n \geq 3$ is even then 
\[
  v_2\left( d_n \right) = n - 2
\]
\end{lemma}
\begin{proof}
  By \cite{MARTIN2004213} 
  \begin{align*}
    d_n = & 
    \frac{
      \left( -1 \right)^{\binom{n}{2}}\left( n^n - 2^n \left( n - 1 \right)^{ n - 1 } \right)
    }{
      \left( n - 2 \right)^2
    } = \\
    & \left( -1 \right)^{\binom{n}{2}}2^{n} \cdot
    \frac{
      \left( \frac{n}{2} \right)^n - \left( n - 1 \right)^{ n - 1 }
    }{
      \left( n - 2 \right)^2
    }
  \end{align*}
  and hence if we denote $m = v_2\left( n - 2 \right)$ we get that 
  \[
    v_2\left( d_n \right) = n - 2m + v_2\left( \left( \frac{n}{2} \right)^n - \left( n - 1 \right)^{ n - 1 } \right)
  \]
Therefore, it suffices to show that 
\[
  v_2\left( \left( \frac{n}{2} \right)^n - \left( n - 1 \right)^{ n - 1 } \right) = 
  2m - n + n - 2 = 2m - 2
\]
Assume that $m \geq 2$
Indeed,
\[
  n - 2 \equiv \epsilon 2^m \mod 2^{2m-1}
\]
for some odd $\epsilon$, and hence 
\[
  n \equiv \epsilon 2^m + 2 \mod 2^{2m-1}
\]
so 
\begin{align*}
  & \left( \frac{n}{2} \right)^n - \left( n - 1 \right)^{n - 1} \equiv \\
  & \left( \epsilon 2^{m - 1} + 1 \right)^n - \left( \epsilon 2^m + 1 \right)^{ n - 1 } \equiv \\
  & \binom{n}{2} \epsilon^2 2^{2m - 2} + n\epsilon 2^{m - 1} + 1 - \left( n - 1 \right) \epsilon 2^m - 1 \equiv 
  \justification{3m - 3 \geq 2m - 1} \\ 
  & \binom{n}{2} \epsilon^2 2^{2m - 2} + 
  \epsilon 2^m \left(  
  \frac{n}{2} - \left( n - 1 \right)
  \right) \equiv \\ 
  & \binom{n}{2} \epsilon^2 2^{2m - 2} + 
  \epsilon 2^m \left(  
    \epsilon 2^{m - 1} + 1 - \epsilon 2^m - 1
  \right) \equiv \\
  & \binom{n}{2} \epsilon^2 2^{2m - 2} + \epsilon 2^m \left( -\epsilon 2^{m-1} \right) \equiv \\
  & \binom{n}{2} \epsilon^2 2^{2m - 2} - \epsilon^2 2^{2m - 1} \equiv \\
  & \left( \binom{n}{2} - 2 \right) \epsilon^2 2^{2m - 2} \mod 2^{2m - 1}
\end{align*}
and in particular its $v_2$ is $2m - 2$ as required.

  Assume finally now that $m = 1$, then 
  \[
    v_2\left( \frac{n}{2} \right) \geq 1
  \]
  and hence 
  \[
    v_2\left( \left( \frac{n}{2} \right)^n \right) \geq n > 0 = v_2\left( \left( n - 1 \right)^{ n - 1 } \right)
  \]
  and hence 
  \[
    v_2\left( \left( \frac{n}{2} \right)^n - \left( n - 1 \right)^{ n - 1 } \right) = 0 = 2m - 2
  \]
  
  Assume  that $m = 2$, then

  \[
    v_2\left( n^n \right) \geq 2n > v_2\left( 2^n \left(  n - 1 \right)^{ n - 1 } \right)
  \]
  and hence 
  \[
    v_2\left( n^n - 2^n \left( n - 1 \right)^{ n - 1 } \right) = n 
  \]
  and hence 
\[
  v_2\left( d_n \right) = n - 2
\]
\end{proof}

\begin{corollary}
  There is an odd prime that divides \( d_n \) for any \( n \geq 4 \).
\end{corollary}
\begin{proof}
  Sketch: The maximal odd divisor of \( d_n \) is 
  \begin{align*}
    & \frac{
      n^n - 2^n \left( n - 1 \right)^{ n - 1 }
    }{
      2^{v_2\left( n^n - 2^n \left( n - 1 \right)^{ n - 1 } \right)} \left( n - 2 \right)^2
    } = \\
    & \frac{
      n^n - 2^n \left( n - 1 \right)^{ n - 1 }
    }{
      2^{n - 2} \left( n - 2 \right)^2
    } = \\
    & \frac{
      \left( \frac{n}{2} \right)^n - \left( n - 1 \right)^{ n - 1 }
    }{
      \left( \frac{n - 2}{2} \right)^2
    }
  \end{align*}
  and the numerator clearly clearly grows faster than the denominator.
\end{proof}

\subsection{Ramification at $2$}
\begin{lemma}
  Let $n \geq 4$ be even.
  Then $L_n$ ramifies at $2$.
\end{lemma}
\begin{proof}
  Let $v$ be a valuation extending $v_2$ to $L_n$ and let $\widehat{L}$ be the completion of $L_n$ with respect to $v$.
  It suffices to show that 
\end{proof}



% Maybe a better proof 
% \begin{proof}
%   \begin{align*}
%     & \frac{
%       d_n
%     }{
%       2^{n-2}
%     } = \\
%     & \frac{
%       \left( -1 \right)^{\binom{n}{2}}\left( n^n - 2^n \left( n - 1 \right)^{ n - 1 } \right)
%     }{
%       2^{n-2} \left( n - 2 \right)^2
%     } = \\
%     & \frac{
%       \left( -1 \right)^{\binom{n}{2}} \left( \left( \frac{n}{2} \right)^n - \left( n - 1 \right)^{ n - 1 } \right)
%     }{
%       \left( \frac{n - 2}{2} \right)^2
%     }
%   \end{align*}
  
% \end{proof}


% \subsection{Reduction to primes that divide $n - 1$ -- Failed Attempt -- $n = 22$ is a counter example}
% We will basically repeat Matin's argument from \cite{MARTIN2004213} and try to generalize it from the case $p = n - 1$ to the case $p \mid n - 1$.
% Assume then that $p \mid n - 1$ is a prime. 
% \begin{lemma}
% \[
%   x^{n-1} \overline{f}_n\left( 2 - \frac{1}{x} \right) = \overline{f}_{n}\left( x \right)
% \]
% where $\overline{f}_n$ is the reduction of $f_n$ modulo $p$.
% \end{lemma}
% \begin{proof}
%   First, 
%   \begin{align*}
%     & \overline{f}_n\left( x \right) = \\
%     & x^{n-1} - 
%     \sum_{i = 1}^{p - 1} \frac{
%       x^{n-1} - 1
%     }{
%       x^{\frac{n - 1}{p}} - 1
%     } = \\ 
%     & x^{n-1} - \sum_{i = 1}^{p - 1} \frac{
%       \left( x^{\frac{n - 1}{p}} \right)^p - 1
%     }{
%       x^{\frac{n - 1}{p}} - 1
%     } = \\
%     & x^{n-1} - \sum_{i = 1}^{p - 1} \frac{
%       \left( x^{\frac{n - 1}{p}} - 1 \right)^p
%     }{
%       x^{\frac{n - 1}{p}} - 1
%     } = \\
%     & x^{n-1} - \sum_{i = 1}^{p - 1} \left( x^{\frac{n - 1}{p}} - 1 \right)^{p - 1}
%   \end{align*}
%   and hence 
%   \begin{align*}
%     & x^{n - 1} \overline{f}_n\left( 2 - \frac{1}{x} \right) = \\
%     & x^{n - 1} \left( \left( 2 - \frac{1}{x} \right)^{n-1} - \sum_{i = 1}^{p - 1} \left( \left( 2 - \frac{1}{x} \right)^{\frac{n - 1}{p}} - 1 \right)^{p - 1} \right) = \\
%     & \left( 2x - 1 \right)^{n - 1} - \sum_{i = 1}^{p - 1} \left( \left( 2x - 1 \right)^{\frac{n - 1}{p}} - x^{\frac{n - 1}{p}} \right)^{p - 1} = \\
%   \end{align*}
  
% \end{proof}
% \begin{lemma}
%   The second Martin's lemma.
% \end{lemma}
% \begin{proof}
%   \[
%     x = 2 - \frac{1}{x^{n - 1}} = 
%     2 - \left( \frac{1}{x^{\frac{n - 1}{p}}} \right)^p = 
%     \left( 2 - \frac{1}{x^{\frac{n - 1}{p}}} \right)^p
%   \]
%   and 
%   \begin{align*}
%     & 2 - \frac{1}{x^{\frac{n - 1}{p}}} = \\
%     & 2 - \frac{
%       -\frac{1}{x} + 1
%     }{
%       x^{\frac{n - 1}{p}} - x^{\frac{n - 1}{p} - 1} 
%     } = \\
%     & 2 + \frac{}{}
%   \end{align*}
  
% \end{proof}

\end{document}